% 按用户真矢量与文楷Regular要求重建；三层职责不是算法流程。
\begin{tikzpicture}[figure picture,foundation picture,
 font=\figurefont\mdseries\fontsize{9}{11}\selectfont,
 trust row/.style={draw=FigureStroke,line width=1pt,rounded corners=1.5mm},
 trust label/.style={align=left,anchor=west,inner sep=0pt},
 trust icon/.style={draw=FigureStroke,line width=.85pt,line cap=round,line join=round}]
 \path[use as bounding box] (0,-2) rectangle (169,64);
 \node at (20,59) {层次};
 \node at (77,59) {关键问题};
 \node at (141.5,59) {责任主体};
 \foreach \y/\fillcol/\stage in {44/FigureBlueFill/属性与规范,26/FigureRedFill/证据能力,8/FigureYellowFill/失效与处置}{
  \draw[trust row,fill=\fillcol] (0,{\y-7}) rectangle (40,{\y+7});
  \draw[trust row,fill=\fillcol] (49,{\y-7}) rectangle (105,{\y+7});
  \draw[trust row,fill=\fillcol] (114,{\y-7}) rectangle (169,{\y+7});
  \node at (20,\y) {\stage};
  \draw[draw=FigureStroke,line width=.9pt] (40,\y)--(49,\y) (105,\y)--(114,\y);
  % 责任记录小图标，不以勾号暗示已经批准。
  \draw[trust icon,fill=white] (118,{\y-3.4})--(123.6,{\y-3.4})--(123.6,{\y+1.6})--(121.8,{\y+3.4})--(118,{\y+3.4})--cycle;
  \draw[trust icon] (121.8,{\y+3.4})--(121.8,{\y+1.6})--(123.6,{\y+1.6});
  \draw[trust icon] (119,{\y+.4})--(122.3,{\y+.4}) (119,{\y-1})--(122.3,{\y-1}) (119,{\y-2.3})--(121.2,{\y-2.3});
 }
 % 规范文件、证据核查、生命周期时钟，对应三个不同问题。
 \draw[trust icon,fill=white] (53,40.6)--(58.6,40.6)--(58.6,45.6)--(56.8,47.4)--(53,47.4)--cycle;
 \draw[trust icon] (56.8,47.4)--(56.8,45.6)--(58.6,45.6);
 \draw[trust icon] (54,44.4)--(57.3,44.4) (54,43)--(57.3,43) (54,41.7)--(56.2,41.7);
 \draw[trust icon] (55.5,27) circle (2.8);
 \draw[trust icon] (57.5,25)--(60,22.5);
 \draw[trust icon] (56.1,8) circle (3.2);
 \draw[trust icon] (56.1,10.2)--(56.1,8)--(58,8);
 \node[trust label] at (63,44) {必须满足什么？\\对谁、什么范围？};
 \node[trust label] at (63,26) {支持多强结论？\\还缺什么证据？};
 \node[trust label] at (63,8) {谁使主张失效？\\何时限制、回退？};
 \node[trust label] at (128,44) {任务责任方\\受影响者、领域规则};
 \node[trust label] at (128,26) {技术人员：论证\\复核者：边界};
 \node[trust label] at (128,8) {有权主体：决策\\运行责任方：处置};
\end{tikzpicture}%
